\section{Full-record lower bounds and tail essentiality}

The lower bounds complement the attainment guarantee in \cref{obj:thm:original-record-attainable-rate} by constructing null and alternative mixtures that remain close when the entire original record is observed. Two constructions supply the required comparisons: paired tents with constant propensity, and coupled continuous coefficient fields in the rough phase. Their separate bounded versions provide witnesses at moment order two.

We use total variation in the convention
\[
\operatorname{TV}(P,Q)=\sup_{E\ \mathrm{measurable}}|P(E)-Q(E)|.
\]
For probability laws \(P\ll Q\), the directed chi-square divergence is
\[
\chi^2(P\Vert Q)
=\int\left(\frac{dP}{dQ}-1\right)^2\,dQ
=\int\left(\frac{dP}{dQ}\right)^2\,dQ-1.
\]
These divergences compare probability laws on the same measurable space.

A finite prior draws one original-record law and then generates the independent sample from that law. Specifically, for
\[
\pi=\sum_{j=1}^{k}\omega_j\delta_{P_j},
\qquad
k\ge1,\quad \omega_j\ge0,\quad \sum_{j=1}^{k}\omega_j=1,
\]
write
\[
\operatorname{supp}_+(\pi)=\{P_j:\omega_j>0\},
\qquad
\mathbb P_{\pi,n}=\sum_{j=1}^{k}\omega_jP_j^{\otimes n}.
\]
Thus \(\mathbb P_{\pi,n}\) mixes product laws after the common latent law has been drawn. Repeated support laws have their weights combined. Normalization ensures that the positive-weight support is nonempty.

For the paired-tent comparison, let \(\lambda\) be Lebesgue probability measure on \([0,1]\), and let \(Q_{a,P}(dy\mid x)\) be the conditional outcome probability kernel in treatment arm \(a\). The following result supplies both rare-mark witnesses under a finite conditional moment and signed-binary witnesses in the matched bounded model.

\begin{lemmav}{lem:paired-tent-full-record-lower}[Paired-tent full-record lower bounds]
Let \(\lambda\) denote Lebesgue probability measure on \([0,1]\), and let \(Q_{a,P}(dy\mid x)\) denote the conditional outcome probability kernel under law \(P\) in treatment arm \(a\in\{0,1\}\). Define the heterogeneity distance, its model suprema, and the fixed witness distance by
\[
d(P)=
\left\{\int_0^1
\left(\tau_P(x)-\int_0^1\tau_P(z)\,\lambda(dz)\right)^2
\,\lambda(dx)\right\}^{1/2},
\qquad
D_v=\sup_{P\in\mathcal M_v}d(P),
\qquad
D_w^{\mathrm b}=\sup_{P\in\mathcal M^{\mathrm b}_w}d(P),
\qquad
d_0=\frac{1}{8\sqrt2},
\]
where the models are defined in \cref{obj:def:model,obj:def:bounded-model}.

Write a normalized finite probability prior as
\[
\pi=\sum_{j=1}^{k}\omega_j\delta_{P_j},
\qquad
\omega_j\ge0,
\qquad
\sum_{j=1}^{k}\omega_j=1.
\]
Define its positive-weight support and full original-record mixture by
\[
\operatorname{supp}_+(\pi)=\{P_j:\omega_j>0\},
\qquad
\mathbb P_{\pi,n}=\sum_{j=1}^{k}\omega_jP_j^{\otimes n}.
\]
For probability laws on the same measurable space, define total variation by
\[
\operatorname{TV}(P,Q)=
\sup_{E\ \mathrm{measurable}}|P(E)-Q(E)|.
\]

For every admissible moment and smoothness tuple \(v=(p,\alpha,\beta,\gamma)\) as in \cref{obj:def:model}, put
\[
q=\frac{p-1}{p},
\qquad
E_0=\frac{2\gamma q}{2\gamma+q},
\qquad
c^{\mathrm e}_v=\frac{4^{-\gamma}}{16\sqrt3}.
\]
There exists a sequence of mark magnitudes \(L:\mathbb N\to\mathbb R\) with
\[
\lim_{n\to\infty}L(n)=+\infty
\]
such that, for every integer \(n\ge2\), there exist a null law \(P_0\) and a normalized finite probability prior \(\pi_1\), with at least one strictly positive weight, satisfying:
\begin{itemize}
\item \textbf{(Null and propensity.)}
The law \(P_0\) belongs to \(H_0(v)\) from \cref{obj:def:null}, and
\[
e_{P_0}(x)=\frac12
\qquad\text{for every }x\in[0,1].
\]

\item \textbf{(Arm support and moments.)}
The magnitude \(L(n)\) is positive. For \(P=P_0\) and every \(P\in\operatorname{supp}_+(\pi_1)\), each arm \(a\in\{0,1\}\) satisfies, for \(\lambda\)-almost every \(x\),
\[
Q_{a,P}\bigl(\{0,-L(n),L(n)\}\mid x\bigr)=1,
\qquad
\int_{\mathbb R}|y|^p\,Q_{a,P}(dy\mid x)\le1.
\]

\item \textbf{(Alternative placement.)}
Every \(P\in\operatorname{supp}_+(\pi_1)\) belongs to \(\mathcal M_v\) from \cref{obj:def:model}, satisfies \(e_P(x)=1/2\) for every \(x\in[0,1]\), and obeys
\[
c^{\mathrm e}_v n^{-E_0}\le d(P)<d_0.
\]
Moreover, \(d_0\le D_v\).

\item \textbf{(Full-record comparison.)}
Define the alternative mixture by
\[
\overline P_1^{(n)}=\mathbb P_{\pi_1,n}.
\]
The same witnesses satisfy
\[
\operatorname{TV}\bigl(P_0^{\otimes n},\overline P_1^{(n)}\bigr)<\frac12.
\]
\end{itemize}
For every such \(n\), the supplied-propensity and original-record minimax testing risks in \cref{obj:def:oracle-risk,obj:def:testing-risk} satisfy
\[
B_n^{\mathrm e}\bigl(v,c^{\mathrm e}_v n^{-E_0}\bigr)\ge\frac25,
\qquad
B_n\bigl(v,c^{\mathrm e}_v n^{-E_0}\bigr)\ge\frac25.
\]

For every admissible smoothness tuple \(w=(\alpha,\beta,\gamma)\) as in \cref{obj:def:bounded-model}, put
\[
c_w^{\mathrm b}=\frac{4^{-\gamma}}{16\sqrt3}.
\]
For every integer \(n\ge2\), there exist a signed-binary null law \(P_0^{\mathrm b}\) and a normalized finite probability prior \(\pi_1^{\mathrm b}\), with at least one strictly positive weight, satisfying:
\begin{itemize}
\item \textbf{(Signed-binary null.)}
The law \(P_0^{\mathrm b}\) belongs to \(H_0^{\mathrm{bin}}(w)\) from \cref{obj:def:binary-null}, and
\[
e_{P_0^{\mathrm b}}(x)=\frac12
\qquad\text{for every }x\in[0,1].
\]

\item \textbf{(Signed-binary alternatives.)}
Every \(P\in\operatorname{supp}_+(\pi_1^{\mathrm b})\) belongs to \(\mathcal M^{\mathrm{bin}}_w\) from \cref{obj:def:binary-model}, satisfies \(e_P(x)=1/2\) for every \(x\in[0,1]\), and obeys
\[
c_w^{\mathrm b}n^{-2\gamma/(4\gamma+1)}
\le d(P)<d_0.
\]
Moreover, \(d_0\le D_w^{\mathrm b}\).

\item \textbf{(Full-record comparison.)}
Define the signed-binary alternative mixture by
\[
\overline P_{1,\mathrm b}^{(n)}
=\mathbb P_{\pi_1^{\mathrm b},n}.
\]
The same witnesses satisfy
\[
\operatorname{TV}\bigl((P_0^{\mathrm b})^{\otimes n},
\overline P_{1,\mathrm b}^{(n)}\bigr)<\frac12.
\]
\end{itemize}
For every such \(n\), the bounded-model minimax testing risk in \cref{obj:def:bounded-risk} satisfies
\[
B_n^{\mathrm b}\bigl(w,c_w^{\mathrm b}n^{-2\gamma/(4\gamma+1)}\bigr)
\ge\frac25.
\]
\end{lemmav}

The fixed propensity in \cref{obj:lem:paired-tent-full-record-lower} makes the comparison applicable to both observation regimes. Every alternative support law meets the separation requirement, while the mixture remains close to the null in the complete original-record experiment. The growing mark magnitudes coexist with a conditional \(p\)-moment bound of one. The signed-binary package supplies a separate bounded comparison at the moment-order-two exponent.

\subsection{Marked records and continuous coefficient fields}

The rough-phase construction uses a normalized marked table \leanref{sym:Cone}{\ensuremath{\mathsf T}} to specify treatment, outcome sign, and mark occurrence jointly. Keeping its six categories explicit preserves the assignment information and the zero outcomes in the likelihood comparison.

\begin{definitionv}{def:marked-table}[Six-category record table \(\mathsf T\)]
The conditional original-record table is
\[
\begin{aligned}
\mathsf T\bigl(A=(1+\ell)/2,Y=hL\mid X=x\bigr)
&=\frac{\varepsilon}{4}
\{1+\ell\xi(x)+h\upsilon(x)+\ell h\zeta(x)\},
&&(\ell,h)\in\{-1,1\}^2,\\
\mathsf T\bigl(A=(1+\ell)/2,Y=0\mid X=x\bigr)
&=\frac{1-\varepsilon}{2}\{1+\ell\xi(x)\},
&&\ell\in\{-1,1\}.
\end{aligned}
\]
Here \(\varepsilon\) is the rare-mark probability, \(L>0\) is the mark magnitude, and \(\xi,\upsilon,\zeta\) may be Borel functions of the covariate. The table is evaluated on inputs for which every displayed probability is nonnegative. Likelihood products include all six record categories.
\end{definitionv}

In \cref{obj:def:marked-table}, the treatment sign \leanref{sym:treatmentSign}{\ensuremath{\ell}} and marked-outcome sign \leanref{sym:markSign}{\ensuremath{h}} index the four nonzero-outcome categories. The coordinate \leanref{sym:xi}{\ensuremath{\xi}} governs treatment probabilities, \leanref{sym:upsilon}{\ensuremath{\upsilon}} governs marked-outcome signs, and \leanref{sym:zeta}{\ensuremath{\zeta}} governs their interaction. Summing over the categories gives total probability one and nonzero-outcome probability \leanref{sym:epsilon}{\ensuremath{\varepsilon}}; the positive mark magnitude is \leanref{sym:Lmark}{\ensuremath{L}}. The remaining two categories retain the treatment coordinate when the outcome is zero.

Continuous coefficient fields are formed from the finite-overlap frame \leanref{sym:Frame}{\ensuremath{(f_i)_{i=0}^{K}}} below, where \(K\) is the fine rank. Its adjacent-cell support gives a continuous construction on the fine partition.

\begin{definitionv}{def:frame-handle}[Finite-overlap frame \(f_i\)]
The finite-overlap frame at rank \(K=J_{\mathrm f}\), the fine histogram rank, consists of the functions
\[
f_i(x)=
\begin{cases}
\cos\!\bigl(\pi(Kx-i)/2\bigr),
&x\in[i/K,(i+1)/K],\\
\sin\!\bigl(\pi(Kx-(i-1))/2\bigr),
&i>0,\ x\in[(i-1)/K,i/K],\\
0,&\text{otherwise},
\end{cases}
\qquad x\in[0,1],\quad 0\le i\le K.
\]
For a coefficient vector
\[
a=(a_0,\ldots,a_K),
\]
the associated frame field is
\[
x\longmapsto\sum_{i=0}^{K}a_i f_i(x).
\]
\end{definitionv}

On each fine cell, the two active frame coordinates have squared values summing to one. The prior construction uses these squares to interpolate coarse paired tents, while the frame coordinates themselves carry the random propensity and outcome coefficients. The following algorithm specifies both branches through a finite collection of sign draws.

\begin{algorithmv}{def:copula-frame-priors}[Copula-frame priors \(\pi^{K,M}_{\nu;v,a,u,\varepsilon,L}\)]
Inputs are the public tuple \(v\), dyadic ranks \(K=J_{\mathrm f}\) and \(M=J_{\mathrm c}\) satisfying \(M\ge2\) and \(K\ge16M\), amplitudes \(0<a,u\le1/16\), a rare-mark probability \(0<\varepsilon<1\), and a mark magnitude \(L>0\).
\begin{enumerate}
\item Set the coupling amplitude and triangular tent to
\[
\kappa=\frac1{16},
\qquad
b_\triangle(t)=2\min\{t,1-t\},
\quad t\in[0,1].
\]
Draw independent uniform signs
\[
\Pr(\sigma_j=1)=\Pr(\sigma_j=-1)=\frac12,
\qquad 1\le j\le M/2.
\]
Define the opposite paired coarse tents by
\[
g_\sigma(x)=
\begin{cases}
\sigma_j b_\triangle(Mx-(2j-2)),
&x\in[(2j-2)/M,(2j-1)/M],\\
-\sigma_j b_\triangle(Mx-(2j-1)),
&x\in[(2j-1)/M,2j/M],
\end{cases}
\quad 1\le j\le M/2.
\]
Their values at coarse-cell boundaries are zero.
\item On each fine cell define the two nonzero frame coordinates by
\[
f_i(x)=\cos\!\bigl(\pi(Kx-i)/2\bigr),
\qquad
f_{i+1}(x)=\sin\!\bigl(\pi(Kx-i)/2\bigr),
\]
\[
x\in[i/K,(i+1)/K],
\qquad 0\le i<K.
\]
Set the other coordinates to zero and form the squared-frame interpolation
\[
\widetilde g_\sigma(x)
=
\sum_{i=0}^{K}g_\sigma(i/K)f_i(x)^2.
\]
\item For each branch \(\nu\in\{0,1\}\), draw coefficient pairs independently conditional on \(\sigma\), with
\[
\Pr_\nu(\lambda_i=l,\eta_i=h\mid\sigma)
=
\frac{1+\nu\kappa g_\sigma(i/K)lh}{4},
\qquad
l,h\in\{-1,1\},
\quad 0\le i\le K.
\]
\item For each coefficient draw, define the conditional-table coordinates
\[
\xi(x)=a\sum_{i=0}^{K}f_i(x)\lambda_i,
\qquad
\upsilon(x)=u\sum_{i=0}^{K}f_i(x)\eta_i,
\]
\[
t_\nu(x)=
-\frac{\nu au\kappa}{1-a^2}\widetilde g_\sigma(x),
\qquad
\zeta(x)=
\xi(x)\upsilon(x)+t_\nu(x)\bigl(1-\xi(x)^2\bigr).
\]
Let \(P\) be the original-record law with uniform covariate \(X\) and the normalized six-category marked conditional table evaluated at these coordinates, rare-mark probability \(\varepsilon\), and mark magnitude \(L\).
\item Output the finite prior obtained by pushing the coefficient distribution forward to original-record laws:
\[
\pi^{K,M}_{\nu;v,a,u,\varepsilon,L}
=
\mathcal L_\nu(P),
\qquad \nu\in\{0,1\}.
\]
\end{enumerate}
\end{algorithmv}

The copula-frame prior \leanref{sym:copulaPrior}{\ensuremath{\pi^{K,M}_{\nu;v,a,u,\varepsilon,L}}} in \cref{obj:def:copula-frame-priors} changes the dependence between the fine coefficient signs across branches. Its propensity-coordinate amplitude is \leanref{sym:aamp}{\ensuremath{a}}, and its marked-outcome-coordinate amplitude is \leanref{sym:uamp}{\ensuremath{u}}. Each coefficient draw determines one original-record law, so the pushforward is a finite probability prior even when several draws determine the same law.

For the model tuning, write \(b\) for the baseline amplitude, and denote the resulting null and alternative priors by \(\pi_0\) and \(\pi_1\). The next result gives their support placement under the phase and rank conditions, together with a separate bounded tuning.

\begin{lemmav}{lem:copula-frame-legality}[Copula-frame legality and moments]
Consider an admissible tuple \(v=(p,\alpha,\beta,\gamma)\) as stipulated in \cref{obj:def:model}. Put
\[
q=\frac{p-1}{p},\qquad S=\alpha+\beta,\qquad
F(p)=\frac{2\alpha}{p-1}+\frac{\beta}{q}+\frac{S}{2\gamma}.
\]
Suppose the following conditions hold:
\begin{itemize}
\item \textbf{(Phase.)} \(F(p)<1\).
\item \textbf{(Dyadic ranks.)} \(K\) and \(M\) are powers of two, with \(16M\le K\) and \(2\le M\).
\item \textbf{(Rank compatibility.)} \(M\le K^{S/\gamma}\).
\end{itemize}
Define the construction amplitudes by
\[
a=\frac{K^{-\alpha}}{16},\qquad b=\frac{K^{-\beta}}{256}.
\]
Let \(\lambda\) denote Lebesgue probability measure on \([0,1]\), and write \(Q_{a,P}(dy\mid x)\) for the conditional outcome kernel in treatment arm \(a\) under \(P\). Define the population effect distance by
\[
d(P)=
\left(
\int_0^1
\left|\tau_P(x)-\int_0^1\tau_P(z)\,\lambda(dz)\right|^2
\,\lambda(dx)
\right)^{1/2},
\]
and let \(d_0\) be the fixed positive distance supplied by the legal model witness.

Use the copula-frame priors of \cref{obj:def:copula-frame-priors} with
\[
u=\frac1{16},\qquad
\varepsilon=(16b)^{1/q},\qquad L=\varepsilon^{-1/p},
\qquad
\pi_0=\pi^{K,M}_{0;v,a,u,\varepsilon,L},\qquad
\pi_1=\pi^{K,M}_{1;v,a,u,\varepsilon,L}.
\]
For every realization of the \(M/2\) coarse signs and the \(K+1\) fine coefficient-sign pairs, the null-branch law belongs to \(H_0(v)\) of \cref{obj:def:null}, and the alternative-branch law belongs to \(\mathcal M_v\) of \cref{obj:def:model}. Each alternative-branch law satisfies
\[
\tau_P(x)=-\frac{2ab\kappa}{1-a^2}\widetilde g_\sigma(x)
\quad\text{for every }x\in[0,1],
\qquad
\int_0^1\tau_P(x)\,\lambda(dx)=0,
\]
\[
2^{-17}K^{-S}\le d(P)\le 2^{-14}K^{-S},
\qquad d(P)<d_0.
\]
Every law in either branch satisfies
\[
\int_{\mathbb R}|y|^p\,Q_{a,P}(dy\mid x)=1
\quad\text{for every }a\in\{0,1\}\text{ and }x\in[0,1].
\]

For the bounded construction, let \(w=(\alpha,\beta,\gamma)\) be an admissible tuple as stipulated in \cref{obj:def:bounded-model}, set \(v=(2,\alpha,\beta,\gamma)\), and impose the same phase, dyadic-rank, and rank-compatibility conditions with \(p=2\). Retain the displayed amplitudes \(a,b\), and use the priors of \cref{obj:def:copula-frame-priors} with
\[
u=2b,\qquad \varepsilon=\frac12,\qquad L=1,
\qquad
\pi_0=\pi^{K,M}_{0;v,a,u,\varepsilon,L},\qquad
\pi_1=\pi^{K,M}_{1;v,a,u,\varepsilon,L}.
\]
For every realization of the coarse signs and fine coefficient-sign pairs, the null-branch law belongs to \(H_0^{\mathrm b}(w)\) of \cref{obj:def:bounded-null}, and the alternative-branch law belongs to \(\mathcal M^{\mathrm b}_w\) of \cref{obj:def:bounded-model}. Each alternative-branch law satisfies the displayed effect identity, centering identity, and distance bounds. Every law in either bounded branch satisfies
\[
\int_{\mathbb R}|y|^2\,Q_{a,P}(dy\mid x)=\frac12
\quad\text{for every }a\in\{0,1\}\text{ and }x\in[0,1].
\]
\end{lemmav}

The support conclusion in \cref{obj:lem:copula-frame-legality} applies to every coefficient realization. The rank-compatibility condition links the coarse effect resolution to the product of the two fine-scale amplitudes. The displayed effect is centered and has distance of order \(K^{-S}\). In the finite-moment tuning, mark rarity and magnitude yield conditional \(p\)-moment exactly one; in the bounded tuning, unit mark magnitude and mark probability one half yield conditional second moment one half. Both tunings retain the same effect identity and distance bounds.

\subsection{Augmentation and complete-record likelihood budgets}

The component comparison retains the complete record through a common augmentation. Write \(B_i=\mathbf1\{Y_i\ne0\}\) for the mark indicator, and let \(\mathcal D_{\partial}\) disclose coefficient pairs at the coarse-cell boundary nodes. The augmentation consists of the covariates, these indicators, and this boundary disclosure. Its common probability law is denoted by \(\mu\).

Conditional on the augmentation, use \(P_0\) and \(P_1\) for the two branch label laws. These are local conditional laws in the component comparison. Each label contains treatment and an outcome sign, including a retained sign when its mark indicator is zero. The intermediate conditional label law is denoted by \(Q\), so that the intermediate augmented probability law is \(\mu\otimes Q\). Recovering the observed outcome from the sign and mark indicator gives the original-record mixtures.

Conditioning on the augmentation fixes the allocation of covariates to fine and coarse cells and identifies which records carry nonzero marks. Boundary disclosure is part of this conditioning. The first and second conditional likelihood budgets, denoted by \(\mathcal A\) and \(\mathcal B\), bound the comparisons from \(P_0\) to \(Q\) and from \(Q\) to \(P_1\), respectively. The following result averages these budgets over the same augmentation law and gives the resulting complete-record comparison.

\begin{lemmav}{lem:full-record-copula-component-bound}[Full-record copula component bound]
Let \(v=(p,\alpha,\beta,\gamma)\) be any real tuple, and suppose:
\begin{itemize}
\item \textbf{(Sample size and ranks.)} \(n,K,M\) are natural numbers, \(n\ge2\), \(K\) and \(M\) are powers of two, \(M\ge2\), \(16M\le K\), and \(2^{40}n\le K\).
\item \textbf{(Amplitudes.)} \(0<a\le1/16\) and \(0<u\le1/16\).
\item \textbf{(Mark parameters.)} \(0<\varepsilon<1\) and \(L>0\).
\end{itemize}
Write \(\pi_0\) and \(\pi_1\) for the two copula-frame priors
\[
\pi_\nu=\pi^{K,M}_{\nu;v,a,u,\varepsilon,L},
\qquad \nu\in\{0,1\},
\]
and define their complete-record mixtures by
\[
\mathbb P_\nu=\int P^{\otimes n}\,\pi_\nu(dP),
\qquad \nu\in\{0,1\}.
\]
Let \(\lambda\) denote Lebesgue probability measure on \([0,1]\). Total variation, denoted by \(\operatorname{TV}\), is the supremum of the absolute difference of event probabilities. The directed chi-square divergence \(\chi^2\) is the integral of the squared Radon--Nikodym density minus one, for probability laws with the stated absolute continuity.

There exist a common augmentation law \(\mu\), conditional label kernels \(P_0,P_1,Q\), and nonnegative measurable, integrable conditional budgets \(\mathcal A,\mathcal B\) with the following properties. The augmentation coordinate is
\[
\left((X_i)_{i=1}^n,(B_i)_{i=1}^n,\mathcal D_{\partial}\right),
\qquad B_i=\mathbf1\{Y_i\ne0\},
\]
where \(\mathcal D_{\partial}\) records the sampled coefficient pairs
\((\lambda_j,\eta_j)\) exactly at nodes satisfying
\[
jM\bmod K=0,\qquad 0\le j\le K.
\]
Its covariate marginal is \(\lambda^{\otimes n}\), and its flag marginal is
\(\operatorname{Bernoulli}(\varepsilon)^{\otimes n}\).

Each label consists of a binary treatment and a binary outcome sign, retaining the sign also when the mark flag is zero. Under branch \(\nu\), the joint law of the augmentation and these labels induced by the latent copula construction is \(\mu\otimes P_\nu\); the disclosure equals the boundary restriction of the sampled coefficients almost surely. Under each branch, the \(M/2\) coarse signs are independent fair signs, and their joint law with the augmentation is the product of their fair-sign law and \(\mu\).

For each branch, recovering the records by retaining \(X_i\) and the treatment label and setting the outcome to the signed magnitude \(L\) when \(B_i=1\), and to zero when \(B_i=0\), sends \(\mu\otimes P_\nu\) to \(\mathbb P_\nu\). Both augmented branch laws have augmentation marginal \(\mu\).

The intermediate augmented law \(\mu\otimes Q\) is a probability law. For \(\mu\)-almost every augmentation, \(P_0,P_1,Q\) are probability laws on the labels and satisfy
\[
Q\ll P_0,\qquad P_1\ll Q,\qquad
\mathcal A\ge0,\qquad \mathcal B\ge0,
\]
\[
1+\chi^2(Q\Vert P_0)\le \exp(\mathcal A),
\qquad
1+\chi^2(P_1\Vert Q)\le \exp(\mathcal B).
\]
Here all three conditional laws and both budgets are evaluated at the full augmentation triple. Their expectations satisfy
\[
\int \mathcal A\,d\mu
\le \frac{2^{38}a^4u^4\varepsilon^2 n^2}{K},
\qquad
\int \mathcal B\,d\mu
\le \frac{2^{56}a^4u^4\varepsilon^2 n^4}{MK^2}.
\]
If
\[
\frac{2^{38}a^4u^4\varepsilon^2 n^2}{K}
+
\frac{2^{56}a^4u^4\varepsilon^2 n^4}{MK^2}
\le 2^{-16},
\]
then
\[
\operatorname{TV}(\mathbb P_0,\mathbb P_1)<\frac18.
\]
\end{lemmav}

\begin{algorithmv}{def:converse-handle}[Branchwise converse priors \(\pi_{0,n,v},\pi_{1,n,v}\)]
Inputs are the sample size \(n\) and the public tuple \(v=(p,\alpha,\beta,\gamma)\); \(q=(p-1)/p\), \(F(p)\) is the exponent-comparison scalar, \(D_p\) is the rough-confounding exponent denominator, \(H_{\mathrm{fine}}\) is the fine-rank multiplier, and \(N_{\mathrm{low}}(v)\) is the rough-construction threshold.
\begin{enumerate}
\item If \(F(p)<1\) and \(n\ge N_{\mathrm{low}}(v)\), select
\[
K=2^{\left\lceil\log_2\!\left(H_{\mathrm{fine}}n^{2/D_p}\right)\right\rceil},
\qquad
M=2^{\left\lfloor\log_2\!\left(K^{(\alpha+\beta)/\gamma}/16\right)\right\rfloor},
\]
and set
\[
a=\frac{K^{-\alpha}}{16},
\qquad
u=\frac1{16},
\qquad
\varepsilon_{n,v}=16^{-1/q}K^{-\beta/q},
\qquad
L_{n,v}=\varepsilon_{n,v}^{-1/p}.
\]
Select the finite copula-frame priors
\[
\pi_{\nu,n,v}
=
\pi^{K,M}_{\nu;v,a,u,\varepsilon_{n,v},L_{n,v}},
\qquad \nu\in\{0,1\}.
\]
\item Otherwise, set
\[
N=2\left\lceil n^{2q/(2\gamma+q)}\right\rceil,
\qquad
h=\frac1N,
\qquad
\varepsilon_{n,v}=h^{\gamma/q},
\qquad
L_{n,v}=h^{-\gamma/(p-1)}.
\]
Select the opposite paired-tent null and alternative priors at this rank, cell width, rare-mark probability, and mark magnitude, denoted respectively by
\[
\pi_{0,n,v},\qquad \pi_{1,n,v}.
\]
\item Output the selected priors and their complete independent-record mixtures
\[
\bigl(\pi_{0,n,v},\pi_{1,n,v}\bigr),
\qquad
\mathbb P_{\nu,n,v}
=
\int P^{\otimes n}\,\pi_{\nu,n,v}(dP),
\quad \nu\in\{0,1\}.
\]
The selected construction satisfies the stipulated balancing, support-law membership, true null-mixture denominator, complete occupancy and transition accounting, and bounded-scale comparison properties.
\end{enumerate}
\end{algorithmv}

The common augmentation in \cref{obj:lem:full-record-copula-component-bound} includes the disclosed boundary coefficients, and the coarse signs remain independent of it under both branches. All three conditional label laws are normalized probability laws. The first expected budget has sample-size and rank factor \(n^2/K\); the second has factor \(n^4/(MK^2)\). Both include the amplitude and mark factor \(a^4u^4\varepsilon^2\). Their stated sum condition gives total variation below one eighth for the recovered complete-record mixtures, including the treatment labels and zero outcomes.

\subsection{Selected priors and the rough lower bound}

The selected converse construction in \cref{obj:def:converse-handle} now has its ingredients. For \(F(p)<1\) and sample size at least the public lower-construction threshold, it selects the copula-frame priors at its specified ranks and finite-moment tuning. The other branch uses the paired-tent package. This selection combines the rough-phase construction with witnesses available for every integer sample size at least two.

At moment order two, the bounded companion uses the bounded copula tuning in the large-sample rough branch and the signed-binary paired-tent package otherwise. Let \(\mathcal V\) denote the admissible tuple domain stipulated in \cref{obj:def:model}. The following definition fixes the bounded prior pair and its complete-record mixtures.

\begin{definitionv}{synth_1}[Bounded converse priors]
Fix \(v=(2,\alpha,\beta,\gamma)\in\mathcal V\), write \(w=(\alpha,\beta,\gamma)\), and let \(n\ge2\). Use the phase functional and public constants of \cref{obj:def:sharp-frontier-scales}. For \(\nu\in\{0,1\}\), define the bounded prior pair \(\pi^{\mathrm b}_{\nu,n,v}\) as follows.

If \(F(2)<1\) and \(n\ge N_{\mathrm{low}}(v)\), put \(S=\alpha+\beta\), let \(K\) be the least power of two at least \(H_{\mathrm{fine}}n^{2/D_2}\), and set
\[
 M=2^{\lfloor\log_2(K^{S/\gamma}/16)\rfloor},
 \qquad a=\frac{K^{-\alpha}}{16},
 \qquad b=\frac{K^{-\beta}}{256}.
\]
Define
\[
 \pi^{\mathrm b}_{\nu,n,v}
   =\pi^{K,M}_{\nu;v,a,\,2b,\,1/2,\,1},
 \qquad \nu\in\{0,1\},
\]
using the finite coefficient draws and record-law pushforward of \cref{obj:def:copula-frame-priors}. This is the bounded tuning of \cref{obj:lem:copula-frame-legality}: it uses the same ranks and amplitudes \(a,b\) as the rough branch of \cref{obj:def:converse-handle}, with outcome-coordinate amplitude \(u=2b\), mark probability \(\varepsilon=1/2\), and mark magnitude \(L=1\).

If \(F(2)\ge1\), or if \(F(2)<1\) and \(n<N_{\mathrm{low}}(v)\), take the signed-binary paired-tent package of \cref{obj:lem:paired-tent-full-record-lower} at sample size \(n\) and tuple \(w\). Writing \(P_0^{\mathrm b}\) for its selected null law and \(\pi_1^{\mathrm b}\) for its selected finite alternative prior, define
\[
 \pi^{\mathrm b}_{0,n,v}=\delta_{P_0^{\mathrm b}},
 \qquad
 \pi^{\mathrm b}_{1,n,v}=\pi_1^{\mathrm b}.
\]
This supplies the signed-binary counterpart of the paired-tent branch and small-sample fallback in \cref{obj:def:converse-handle}.

In either branch, the associated complete-record mixtures are
\[
 \mathbb P^{\mathrm b}_{\nu,n,v}
   =\int P^{\otimes n}\,\pi^{\mathrm b}_{\nu,n,v}(dP),
 \qquad \nu\in\{0,1\}.
\]
\end{definitionv}

The bounded prior pair \(\pi^{\mathrm b}_{\nu,n,v}\) in \cref{obj:synth_1} is defined for every admissible moment-order-two tuple and every \(n\ge2\). Its complete-record mixture \(\mathbb P^{\mathrm b}_{\nu,n,v}\) uses the same draw-once convention as \(\mathbb P_{\pi,n}\). The next result assembles support placement, likelihood control, and the testing-risk lower bound in the rough phase.

\begin{lemmav}{lem:rough-full-record-sharp-lower}[Sharp rough-phase record lower bound]
Use the selected converse priors and complete-record mixtures of \cref{obj:def:converse-handle}, and the phase functional, exponents, multipliers, and scales of \cref{obj:def:sharp-frontier-scales}. Let \(\lambda\) denote Lebesgue probability measure on \([0,1]\). Define the heterogeneity distance, model supremum, and witness distance by
\[
d(P)=\left(\int_{[0,1]}
\left(\tau_P(x)-\int_{[0,1]}\tau_P(z)\,d\lambda(z)\right)^2
\,d\lambda(x)\right)^{1/2},
\qquad
D_v=\sup_{P\in\mathcal M_v}d(P),
\qquad
d_0=\frac{1}{8\sqrt2}.
\]
For probability laws, \(\operatorname{TV}\) denotes the supremum of the absolute differences between their probabilities of the same measurable event.

Suppose that:
\begin{itemize}
\item \textbf{(Admissibility.)} \(v=(p,\alpha,\beta,\gamma)\) is admissible as in \cref{obj:def:model}.
\item \textbf{(Rough phase.)} \(F(p)<1\).
\item \textbf{(Sample size.)} \(n\) is an integer with \(n\ge2\).
\end{itemize}
The finite priors \(\pi_{0,n,v}\) and \(\pi_{1,n,v}\) have nonnegative weights summing to one. Every positive-weight null support law belongs to \(H_0(v)\) of \cref{obj:def:null}, and every positive-weight alternative support law satisfies
\[
P\in\mathcal M_v,
\qquad
c_vn^{-E_4(p)}\le d(P)<d_0.
\]
Moreover,
\[
d_0\le D_v,
\qquad
\operatorname{TV}(\mathbb P_{0,n,v},\mathbb P_{1,n,v})<\frac12,
\qquad
B_n\!\left(v,c_vn^{-E_4(p)}\right)\ge\frac25,
\]
with testing risk as in \cref{obj:def:testing-risk}.

If \(n\ge N_{\mathrm{low}}(v)\), the selected fine and coarse ranks \(K,M\) of \cref{obj:def:converse-handle} satisfy the rank conditions of \cref{obj:lem:copula-frame-legality}. In particular, they are powers of two and satisfy
\[
16M\le K,
\qquad
2\le M\le K^{S/\gamma},
\qquad
2^{40}n\le K.
\]
With the selected rare-mark probability \(\varepsilon_{n,v}\) of \cref{obj:def:converse-handle}, the deterministic budget bounds are
\[
\begin{aligned}
\frac{2^{38}(K^{-\alpha}/16)^4(1/16)^4
      \varepsilon_{n,v}^{\,2}n^2}{K}
&\le \frac{2^{-10}}{H_{\mathrm{fine}}},\\
\frac{2^{56}(K^{-\alpha}/16)^4(1/16)^4
      \varepsilon_{n,v}^{\,2}n^4}{MK^2}
&\le \frac{2^{13}}{H_{\mathrm{fine}}^2}.
\end{aligned}
\]
The same complete-record mixtures then satisfy
\[
\operatorname{TV}(\mathbb P_{0,n,v},\mathbb P_{1,n,v})<\frac18.
\]

For the bounded experiment, suppose that:
\begin{itemize}
\item \textbf{(Admissibility.)} \(w=(\alpha,\beta,\gamma)\) is admissible as in \cref{obj:def:bounded-model}; write \(v=(2,\alpha,\beta,\gamma)\).
\item \textbf{(Rough phase.)} \(F(2)<1\).
\item \textbf{(Sample size.)} \(n\) is an integer with \(n\ge2\).
\end{itemize}
Let \(\pi^{\mathrm b}_{0,n,v}\) and \(\pi^{\mathrm b}_{1,n,v}\) be the selected bounded converse priors of \cref{obj:synth_1}, and let
\(\mathbb P^{\mathrm b}_{0,n,v}\) and \(\mathbb P^{\mathrm b}_{1,n,v}\) denote their respective complete-record mixtures. Both finite priors have nonnegative weights summing to one. Every positive-weight null support law belongs to \(H_0^{\mathrm b}(w)\) of \cref{obj:def:bounded-null}, and every positive-weight alternative support law satisfies
\[
P\in\mathcal M^{\mathrm b}_w,
\qquad
c_w^{\mathrm b}n^{-E_4(2)}\le d(P)<d_0.
\]
Furthermore,
\[
\operatorname{TV}
\bigl(\mathbb P^{\mathrm b}_{0,n,v},\mathbb P^{\mathrm b}_{1,n,v}\bigr)
<\frac12,
\qquad
B_n^{\mathrm b}\!\left(w,c_w^{\mathrm b}n^{-E_4(2)}\right)\ge\frac25,
\]
with bounded testing risk as in \cref{obj:def:bounded-risk}.
\end{lemmav}

Under \(F(p)<1\), \cref{obj:lem:rough-full-record-sharp-lower} gives the rough separation exponent for every \(n\ge2\). Above the public lower-construction threshold, the selected ranks also meet the component-comparison conditions and give the stronger total-variation bound of one eighth. The deterministic budget bounds control the upper bounds on the expectations of \(\mathcal A\) and \(\mathcal B\) in \cref{obj:lem:full-record-copula-component-bound}. The bounded conclusion uses its own outcome tuning at \(p=2\), preserving the rough exponent and the finite-sample risk guarantee.

\subsection{Tail essentiality and the moment-order-two face}

The final result combines the rough and paired-tent branches with attainment. Its bounded-prior comparison quantifies over normalized finite priors in the convention fixed above: all support restrictions apply to positive-weight laws, and each mixture draws its latent law once for the complete sample. For each fixed admissible tuple with moment order below two, the comparison has a tuple-dependent eventual threshold. On the moment-order-two face, the selected bounded constructions give witnesses for every \(n\ge2\).

\begin{theoremv}{thm:tail-essential-full-record-converse}[Tail essentiality and bounded witnesses]
Fix an admissible tuple \(v=(p,\alpha,\beta,\gamma)\) as in \cref{obj:def:model}, and write \(w=(\alpha,\beta,\gamma)\). Use the scales, constants, thresholds, and total tests of \cref{obj:def:sharp-frontier-scales}, and the selected priors, mark magnitudes, and complete-record mixtures of \cref{obj:def:converse-handle}.

Let \(\lambda\) denote Lebesgue probability measure on \([0,1]\). Define the heterogeneity distance and its model supremum by
\[
d(P)=
\left\{\int_{[0,1]}
\left(\tau_P(x)-\int_{[0,1]}\tau_P(z)\,\lambda(dz)\right)^2
\,\lambda(dx)\right\}^{1/2},
\qquad
D_v=\sup_{P\in\mathcal M_v}d(P).
\]
For probability laws, \(\operatorname{TV}\) denotes the supremum of the absolute differences of their probabilities over measurable events.

A normalized finite prior on original-record laws is
\[
\pi=\sum_{j=1}^{k}\omega_j\delta_{P_j},
\qquad
k\ge1,\quad \omega_j\ge0,\quad \sum_{j=1}^{k}\omega_j=1.
\]
Its positive-weight support and complete-record mixture are defined by
\[
\operatorname{supp}_+(\pi)=\{P_j:\omega_j>0\},
\qquad
\mathbb P_{\pi,n}=\sum_{j=1}^{k}\omega_jP_j^{\otimes n}.
\]

For every integer \(n\ge2\), both \(\pi_{0,n,v}\) and \(\pi_{1,n,v}\) are normalized finite priors, and
\[
\begin{gathered}
\operatorname{supp}_+(\pi_{0,n,v})\subseteq H_0(v),\\
\operatorname{supp}_+(\pi_{1,n,v})
\subseteq
\{P\in\mathcal M_v:c_v\rho_n(v)\le d(P)<D_v\},\\
0<c_v\rho_n(v)<D_v,\\
\operatorname{TV}(\mathbb P_{0,n,v},\mathbb P_{1,n,v})<\frac12,
\qquad
B_n(v,c_v\rho_n(v))\ge\frac25,
\end{gathered}
\]
where the classes and testing risk are those of
\cref{obj:def:null,obj:def:model,obj:def:testing-risk}.
For every positive-weight law in either prior, each conditional outcome arm assigns probability one to
\(\{0,-L_{n,v},L_{n,v}\}\) for \(\lambda\)-almost every covariate value. Moreover,
\[
L_{n,v}\longrightarrow\infty
\qquad(n\longrightarrow\infty).
\]

Use the matched bounded model \(\mathcal M^{\mathrm b}_w\) and null
\(H_0^{\mathrm b}(w)\) of
\cref{obj:def:bounded-model,obj:def:bounded-null}.
If \(p<2\), then
\[
\frac{\rho_n(v)}{\rho_n^{\mathrm b}(w)}
\longrightarrow\infty
\qquad(n\longrightarrow\infty).
\]
For each such fixed \(v\), there exists an integer \(N(v)\ge2\) such that, for every integer \(n\ge N(v)\) and every pair of normalized finite priors \(\pi_0,\pi_1\),
\[
\left.
\begin{gathered}
\operatorname{supp}_+(\pi_0)\subseteq H_0^{\mathrm b}(w),\\
\operatorname{supp}_+(\pi_1)\subseteq
\{P\in\mathcal M^{\mathrm b}_w:c_v\rho_n(v)\le d(P)\}
\end{gathered}
\right\}
\quad\Longrightarrow\quad
\operatorname{TV}(\mathbb P_{\pi_0,n},\mathbb P_{\pi_1,n})
\ge\frac12.
\]

If \(p=2\), the following conclusions hold for every integer \(n\ge2\).
Let \(\pi^{\mathrm b}_{\nu,n,v}\), for \(\nu\in\{0,1\}\), denote the selected bounded witness priors: the bounded rough priors of
\cref{obj:lem:rough-full-record-sharp-lower} when
\(F(2)<1\) and \(n\ge N_{\mathrm{low}}(v)\), and the signed-binary paired-tent priors of
\cref{obj:lem:paired-tent-full-record-lower} otherwise. Define
\[
\mathbb P^{\mathrm b}_{\nu,n,v}
=\mathbb P_{\pi^{\mathrm b}_{\nu,n,v},n},
\qquad \nu\in\{0,1\}.
\]
Both priors are normalized and have nonempty positive-weight supports. Their positive-weight masses satisfy
\[
\sum_{P\in\operatorname{supp}_+(\pi^{\mathrm b}_{\nu,n,v})}
\pi^{\mathrm b}_{\nu,n,v}(\{P\})=1,
\qquad \nu\in\{0,1\}.
\]
Let \(d_0>0\) be the fixed witness distance of
\cref{obj:lem:causal-nonempty}. The support, total-variation, and risk conclusions are
\[
\begin{gathered}
\operatorname{supp}_+(\pi^{\mathrm b}_{0,n,v})
\subseteq H_0^{\mathrm b}(w),
\qquad
\operatorname{supp}_+(\pi^{\mathrm b}_{0,n,v})
\subseteq H_0(v),\\
\operatorname{supp}_+(\pi^{\mathrm b}_{1,n,v})
\subseteq
\{P\in\mathcal M^{\mathrm b}_w:c_v\rho_n(v)\le d(P)<d_0\},\\
\operatorname{supp}_+(\pi^{\mathrm b}_{1,n,v})
\subseteq
\{P\in\mathcal M_v:c_v\rho_n(v)\le d(P)\},\\
\operatorname{TV}(\mathbb P^{\mathrm b}_{0,n,v},
                  \mathbb P^{\mathrm b}_{1,n,v})<\frac12,\\
\frac25\le B_n^{\mathrm b}(w,c_v\rho_n(v))
\le B_n(v,c_v\rho_n(v)),
\end{gathered}
\]
with bounded risk defined in \cref{obj:def:bounded-risk}.

Define the bounded-model distance supremum by
\[
D_w^{\mathrm b}=\sup_{P\in\mathcal M^{\mathrm b}_w}d(P).
\]
On this face, the scales, multipliers, thresholds, and distance suprema agree:
\[
\begin{gathered}
\rho_n(v)=\rho_n^{\mathrm b}(w),
\qquad c_v=c_w^{\mathrm b},
\qquad C_v=C_w^{\mathrm b},\\
N_v=N_w^{\mathrm b},
\qquad D_v=D_w^{\mathrm b}.
\end{gathered}
\]
The capped critical radii of
\cref{obj:def:critical-radius,obj:def:bounded-radius,obj:def:binary-radius}
satisfy
\[
\begin{gathered}
c_v\rho_n(v)\le r_n^*(v)\le C_v\rho_n(v),\\
c_w^{\mathrm b}\rho_n^{\mathrm b}(w)
\le r_n^{*,\mathrm b}(w)=r_n^{*,\mathrm{bin}}(w)
\le C_w^{\mathrm b}\rho_n^{\mathrm b}(w).
\end{gathered}
\]

Let \(\mathsf R_n(P,\phi)\) denote the expected rejection probability under \(n\) independent original records and the public randomization of
\cref{obj:def:original-record-experiment}. The total tests satisfy, for every law in the indicated null class,
\[
\begin{aligned}
\mathsf R_n(P,\phi^*_{n,v})&\le0.0075
&&\quad(P\in H_0(v)),\\
\mathsf R_n(P,\phi^{*,\mathrm b}_{n,w})&\le0.0075
&&\quad(P\in H_0^{\mathrm b}(w)).
\end{aligned}
\]
Whenever \(C_v\rho_n(v)<D_v\), both separated alternative classes are nonempty:
\[
\begin{gathered}
\{P\in\mathcal M_v:C_v\rho_n(v)\le d(P)\}\ne\varnothing,\\
\{P\in\mathcal M^{\mathrm b}_w:
C_w^{\mathrm b}\rho_n^{\mathrm b}(w)\le d(P)\}\ne\varnothing.
\end{gathered}
\]
Under this same condition, every law in the corresponding separated class satisfies
\[
\begin{aligned}
1-\mathsf R_n(P,\phi^*_{n,v})&\le0.0075
&&\quad(P\in\mathcal M_v,\ C_v\rho_n(v)\le d(P)),\\
1-\mathsf R_n(P,\phi^{*,\mathrm b}_{n,w})&\le0.0075
&&\quad(P\in\mathcal M^{\mathrm b}_w,\\
C_w^{\mathrm b}\rho_n^{\mathrm b}(w)\le d(P)).
\end{aligned}
\]
For \(n\ge N_v\), the common upper separation satisfies
\[
\begin{gathered}
C_v\rho_n(v)\le\frac{d_0}{2},\\
C_v\rho_n(v)<D_v,
\qquad
C_w^{\mathrm b}\rho_n^{\mathrm b}(w)<D_w^{\mathrm b}.
\end{gathered}
\]
At \(n=2,3\), both total tests have zero rejection probability at every sample:
\[
\phi^*_{n,v}(z)=0,
\qquad
\phi^{*,\mathrm b}_{n,w}(z)=0
\qquad\text{for every }z.
\]
\end{theoremv}

The distinction in \cref{obj:thm:tail-essential-full-record-converse} concerns the outcome distributions that sustain close full-record comparisons at the finite-moment separation scale. The selected finite-moment witnesses have growing mark magnitudes and total variation below one half for every \(n\ge2\). For each fixed tuple with \(p<2\), the bounded scale is asymptotically smaller, and every normalized finite bounded-prior pair satisfying the stated support restrictions has total variation at least one half once \(n\ge N(v)\). This threshold depends on the fixed tuple.

At \(p=2\), the bounded copula construction and the signed-binary paired-tent construction together supply normalized, nonempty bounded witnesses for every sample size. The common scales and constants then connect these lower comparisons to the attained upper bounds. The power conclusions apply when the common upper separation is below the model distance supremum, a condition guaranteed by the public threshold \(N_v\); the capped-radius conclusions retain the conventions of \cref{obj:def:critical-radius,obj:def:bounded-radius,obj:def:binary-radius}.
